% Part: many-valued-logic
% Chapter: syntax-and-semantics
% Section: formulas

\documentclass[../../../include/open-logic-section]{subfiles}

\begin{document}

\olfileid{mvl}{syn}{fml}

\olsection{\usetoken{P}{formula}}

\begin{defn}[فارمول]
\ollabel{defn:formulas}
د بياني ژبې~$\Lang L$ د \emph{!!{formula}s} سټ~$\Frm[L]$
په استقرايي ډول داسې تعريفېږي:
\begin{enumerate}
\item هر !!{propositional variable}~$\Obj p_i$ يو اټومي
  !!{formula} دے۔
\item د~$\Lang L$ هر $0$-ځايه نښلوونکے (بياني ثابت)
  يو اټومي !!{formula} دے۔
\item که~$\star$ د~$\Lang L$ يو $n$-ځايه نښلوونکے وي، او
  $!A_1$، \dots، $!A_n$ !!{formula}s وي، نو
  $\star(!A_1, \dots, !A_n)$ هم !!a{formula} دے۔
\tagitem{limitClause}{له دې پرته هېڅ څيز !!a{formula} نۀ دے۔}{}
\end{enumerate}
که~$\star$ يو $1$-ځايه نښلوونکے وي، نو
$\star(!A_1)$ به ډېر ځله په ساده ډول $\star !A_1$ ليکو۔
که~$\star$ يو $2$-ځايه نښلوونکے وي، نو
$\star(!A_1,!A_2)$ به ډېر ځله د $(!A_1 \star !A_2)$
په بڼه ليکو۔
\end{defn}

د معمول په څېر، ډېر ځله به تر ټولو بېروني قوسونه له ليکلو پرېږدو۔

\begin{ex}
  په معياري ژبه~$\Lang{L_0}$ کښې
  $\Obj p_1 \lif (\Obj p_1 \land \lnot \Obj p_2)$ يو فارمول دے۔
  د حاصل‌ضرب منطق په ژبه کښې به د دې پر ځاے
  $\Obj p_1 \lif (\Obj p_1 \odot \lnot \Obj p_2)$ وليکو۔
  که ژبې ته يو $1$-ځايه~$\triangle$ هم ور زيات کړو، نو
  $\triangle (\Obj p_1 \land \Obj p_2) \lif (\triangle \Obj p_1 \land \triangle \Obj p_2)$
  غوندې فارمولونه به هم ولرو۔
\end{ex}

\end{document}
